% Regression: p3_slidecd — imported current-behavior fixture from the handoff corpus.
% Formula: e = <(3,2),(3,3)> of a 5x5 clip (K = L = 2) reads as a physical
% p3_slidecd — arXiv:1804.04964 (normal PEPS FT), Corollary 8 sketch,
% p. 24 (top figure).
% Sliding windows joined by arrows: a virtual operation on the bond
% e = <(3,2),(3,3)> of a 5x5 clip (K = L = 2) reads as a physical
% operation on any of four 2x2 windows; consecutive windows differ by
% one lattice step (west, west, then south) and the paper joins the
% panels by arrows.
%
% Window sequence (the edge e is fixed in every panel):
%   W1 = (2-3, 3-4)  ->  W2 = (2-3, 2-3)  ->  W3 = (2-3, 1-2)
%   ->  W4 = (3-4, 1-2).
%
% Hue mapping: paper red window -> slot=selected (Action) outline; paper
% fat green edge -> distinguished edge (Action).  The window/edge
% contrast survives by line weight only (0.80pt outline vs the
% emphasis-width edge) — noted.
\documentclass[varwidth=32cm, border=6pt]{standalone}
\usepackage{amsmath, amssymb}
\usepackage{tenkz}
\begin{document}
\begin{tenkzcd}[column sep=large]
  \begin{tenkzlattice}[rows=5, cols=5, boundary legs, compact]
    \tnregion[slot=selected, outline]{(2-3,3-4)}
    \tnedge{(3,2)-(3,3)}
  \end{tenkzlattice}
  \arrow[r]
  &
  \begin{tenkzlattice}[rows=5, cols=5, boundary legs, compact]
    \tnregion[slot=selected, outline]{(2-3,2-3)}
    \tnedge{(3,2)-(3,3)}
  \end{tenkzlattice}
  \arrow[r]
  &
  \begin{tenkzlattice}[rows=5, cols=5, boundary legs, compact]
    \tnregion[slot=selected, outline]{(2-3,1-2)}
    \tnedge{(3,2)-(3,3)}
  \end{tenkzlattice}
  \arrow[r]
  &
  \begin{tenkzlattice}[rows=5, cols=5, boundary legs, compact]
    \tnregion[slot=selected, outline]{(3-4,1-2)}
    \tnedge{(3,2)-(3,3)}
  \end{tenkzlattice}
\end{tenkzcd}
\end{document}
